\documentclass[11pt]{article}
\usepackage[margin=1in]{geometry}
\usepackage[T1]{fontenc}
\usepackage{lmodern,microtype,hyperref,enumitem,amsmath,amssymb,xcolor,booktabs}
\usepackage[numbers,sort&compress]{natbib}
\hypersetup{colorlinks=true,urlcolor=blue,citecolor=blue,linkcolor=blue}
\title{Q15: Comprehension Check}
\author{DLS 2026 --- Kukla}
\date{October 2026}
\begin{document}
\maketitle
\textbf{Comprehension check.} A team wants their assistant to answer questions over a 10-million-token
internal wiki that updates daily. They propose training the facts into the model's weights. Give the main
drawback of that approach and the standard alternative, and state one reason the alternative also reduces
hallucination.

\vspace{0.8cm}\hrule\vspace{4pt}
{\small\textbf{Answer.} Baking a daily-updating 10M-token corpus into weights is impractical: it far exceeds
any context window to train on at once, and the facts go stale immediately---every daily update would
require retraining/fine-tuning. The standard alternative is \textbf{retrieval-augmented generation}: index
the wiki as vectors and retrieve the relevant chunks per query into the context, so updates are just
re-indexing (no retraining) and the corpus can be arbitrarily large. It also reduces hallucination because
the answer is \emph{grounded} in the retrieved passages and can cite them, rather than relying on the
model's parametric guess.}
\end{document}
